% META: {"id":"TCOMPL-LOG-CO-002","chapitre":"TCOMPL-LOGARITHME-HISTORIQUE","type_objet":"corrige","exercice_ref":"TCOMPL-LOG-EX-002","capacites_codes":["C2"],"sources_inspiration":[],"status":"approved","origin":"generated","status_history":["generated","audited","accepted","approved"],"release_acceptance":"RELEASE_OWNER_BATCH_ACCEPTANCE","acceptance_closure_digest":"sha256:9b3ccf9a81c5520fbb7e03b7b2d3e7bf2057a02834b6d908bf3be5f49f3b553f","acceptance_receipt_digest":"sha256:4fad86f0282e0fa4e0419cca1ad6379ae7af5a280c04a4a90880f90a1c044242"}
% BEGIN-VERIFY
% from sympy import symbols, ln, solve, Eq
% x = symbols('x')
% sols = solve(Eq(ln(2*x+1), ln(x) + ln(3)), x)
% assert sols == [1]
% END-VERIFY
\begin{corrige}{TCOMPL-LOG-EX-002}
\textbf{Domaine.} Il faut $2x+1>0$ (soit $x>-\frac{1}{2}$) et $x>0$ : le domaine est $]0,+\infty[$.

\textbf{Résolution.} Sur ce domaine, $\ln(x)+\ln(3) = \ln(3x)$ (équation fonctionnelle). L'équation devient :
\[
  \ln(2x+1) = \ln(3x) \iff 2x+1 = 3x \iff x = 1,
\]
car $\ln$ est strictement croissante donc injective. On vérifie que $x=1>0$ : $\boxed{S=\{1\}}$.
\end{corrige}
